\documentclass[10pt]{article}
\usepackage[T1]{fontenc}
\usepackage[utf8]{inputenc}
\usepackage[french]{babel}
\usepackage{amsmath,amssymb,mathrsfs,mathpazo}
\usepackage{geometry}
\geometry{a4paper,margin=2cm}
\usepackage{graphicx,array,tabularx,tabularray,colortbl,multirow,multicol,booktabs,enumitem,pifont,ulem,fancybox,diagbox,cancel}
\usepackage{pstricks,pstricks-add,pst-plot,pst-tree,pst-node,pst-text,pst-func,pst-math,pst-bspline,pst-3dplot,pst-eucl,pst-solides3d}
\let\pscalculate\undefined
\usepackage{pst-bezier}
\usepackage{tikz,pgfplots,tkz-tab}
\usetikzlibrary{arrows,arrows.meta,patterns,calc,positioning,trees,decorations.markings,shapes.geometric,shapes.misc,shapes.arrows,intersections,decorations.pathmorphing,decorations.pathreplacing,backgrounds,shadows,babel}
\pgfplotsset{compat=1.18}
\usepgfplotslibrary{fillbetween}
\usepackage[np]{numprint}
\usepackage{listings,makecell,pas-tableur,twoopt}
\let\barre\undefined
\usepackage[pstricks]{tablvar}
\definecolor{darkgreen}{RGB}{0,100,0}
\providecommand{\red}{\color{red}}
\providecommand{\blue}{\color{blue}}
\providecommand{\green}{\color{green}}
\providecommand{\black}{\color{black}}
\providecommand{\white}{\color{white}}
\providecommand{\gray}{\color{gray}}
\providecommand{\cyan}{\color{cyan}}
\providecommand{\magenta}{\color{magenta}}
\providecommand{\yellow}{\color{yellow}}
\providecommand{\violet}{\color{violet}}
\providecommand{\orange}{\color{orange}}
\providecommand{\euro}{\texteuro}
\providecommand{\EUR}{\texteuro}
\providecommand{\ieme}{\textsuperscript{e}}
\providecommand{\degres}{\ensuremath{^\circ}}
\providecommand{\degree}{\ensuremath{^\circ}}
\providecommand{\vv}[1]{\overrightarrow{#1}}
\providecommand{\Coord}[1]{#1}
\providecommand{\decosix}{}
\providecommand{\starredbullet}{\ensuremath{\bullet}}
\pagestyle{empty}
\setlength{\parindent}{0pt}
\usepackage{adjustbox}\DeclareUnicodeCharacter{2236}{:}\DeclareUnicodeCharacter{2190}{\ensuremath{\leftarrow}}
\providecommand{\R}{\mathbb{R}}
\providecommand{\N}{\mathbb{N}}
\providecommand{\Z}{\mathbb{Z}}
\providecommand{\Q}{\mathbb{Q}}
\providecommand{\D}{\mathbb{D}}
\providecommand{\C}{\mathbb{C}}
\providecommand{\e}{{\mathrm{e}}}
\providecommand{\vect}[1]{\overrightarrow{#1}}
\providecommand{\vectt}[1]{\overrightarrow{\mathrm{#1}}}
\providecommand{\vvt}[1]{\overrightarrow{\mathrm{#1}}}
\providecommand{\np}[2][]{{#2~#1}}
\providecommand{\vv}[1]{{\overrightarrow{#1}}}
\renewcommand{\d}{{\mathrm{d}}}
\providecommand{\degres}{{^\circ}}
\providecommand{\degre}{{^\circ}}
\providecommand{\blue}{\color{blue}}
\providecommand{\red}{\color{red}}
\providecommand{\enskip}{\quad}
\providecommand{\llbracket}{\lbrack\!\lbrack}
\providecommand{\rrbracket}{\rbrack\!\rbrack}
\providecommand{\euro}{{\text{€}}}
\providecommand{\no}{{\text{n°}}}
\providecommand{\interval}[2]{{[#1;#2]}}
\providecommand{\intervFO}[2]{{[#1;#2[}}
\providecommand{\dotfill}{\ldots}

\providecommand{\ds}{\displaystyle}
\providecommand{\textding}[1]{#1}
\providecommand{\Oij}{\(\mathrm{O};\vec{\imath},\vec{\jmath}\)}
\providecommand{\Oijk}{\(\mathrm{O};\vec{\imath},\vec{\jmath},\vec{k}\)}
\providecommand{\Ouv}{\(\mathrm{O};\vec{u},\vec{v}\)}
\providecommand{\intervFF}[2]{\ensuremath{[#1;#2]}}
\renewcommand{\euro}{\ifmmode\text{\texteuro}\else\texteuro\fi}
\DeclareUnicodeCharacter{00B0}{\ifmmode^{\circ}\else\textdegree\fi}
\DeclareUnicodeCharacter{20AC}{\ifmmode\text{\texteuro}\else\texteuro\fi}
% Énoncés : Éducation nationale. Transcription et corrigés : APMEP, Denis Vergès et contributeurs. Découpage et classement : CoopMaths, Sébastien Lozano, Éric Elter et contributeurs. Adaptation de la présentation : Maths Explorer.
% CC BY-SA : https://coopmaths.fr/www/about/
% https://coopmaths.fr/alea/static/e3c/2020/tex/e3c_2020_00_sujet35_4_cor.tex
\begin{document}
\noindent\textit{APMEP / CoopMaths. CC BY-SA. Présentation adaptée par Maths Explorer.}\par\medskip
\subsection*{Partie A : Lecture graphique}

\paragraph{1.} La tangente en \(A\) est horizontale donc \(f'(0) = 0\).

La tangente en \(B\), qui est la droite \((BC)\), a pour coefficient directeur :
\[
f'(2) = \dfrac{6 - 2}{1 - 0} = 4.
\]

\subsection*{Partie B : Calcul algébrique}

\(f(x) = \e^x (2x + 2)\).

\paragraph{2.} En tant que produit de fonctions dérivables sur \(\mathbb{R}\), \(f\) est dérivable sur \(\mathbb{R}\) et :
\[
f'(x) = \e^x (2x + 2) + 2\e^x = \e^x(2x + 2 + 2) = \e^x(2x + 4).
\]

\paragraph{3.} On sait que, quel que soit le réel \(x\), \(\e^x > 0\), donc le signe de \(f'(x)\) est celui de \(2x + 4\).

\[
2x + 4 > 0 \iff x > -2
\]

Donc \(f\) est décroissante sur \(\left]-\infty\,;\,-2\right[\), puis croissante sur \(\left]-2\,;\,+\infty\right[\), avec un minimum en :
\[
f(-2) = \e^{-2}(2 \times (-2) + 2) = -2 \e^{-2} \approx -0{,}271.
\]

\begin{center}
\begin{tikzpicture}
\tkzTabInit[lgt=3, espcl=3]{$x$ / 1, {$\e^x$} / 1, {$2x + 4$} / 1, {Signe de $f'(x)$} / 1, {$f$} / 2}{${-\infty}$, ${-2}$, ${+\infty}$}
\tkzTabLine{,+,,+,}
\tkzTabLine{,-,0,+,}
\tkzTabLine{,-,0,+,}
\tkzTabVar{+/{$$},-/{$-2\e^{-2} \approx -0{,}271$},+/{$$}}{/}
\end{tikzpicture}
\end{center}

\paragraph{4.} On sait que l'équation réduite de la tangente au point d'abscisse \(0\) est :
\[
y - f(0) = f'(0)(x - 0).
\]
Avec \(f(0) = 1 \times 2 = 2\) et \(f'(0) = 1 \times (0 + 4) = 4\), l'équation devient :
\[
y - 2 = 4x \quad \text{ou} \quad y = 4x + 2.
\]

\paragraph{5.}
La tangente en \(A\) est parallèle à l'axe des abscisses.  
\(A\) a pour abscisse \(-2\), donc :
\[
f'(-2) = \e^{-2} \times (2 \times (-2) + 4) = \e^{-2} \times 0 = 0,
\]
la tangente en \(A\) est parallèle à l'axe des abscisses.

L'équation de la tangente étant \(y = 4x + 2\), on a bien pour \(x = 1\) (abscisse de C) :
\[
y = 4 \times 1 + 2 = 6 \text{ (ordonnée de C).}
\]
La tangente en \(B\) contient bien le point \(C\).

\end{document}
